% !TeX root = ../../../progress-sem1.tex
% 4.x การประเมินผลการดำเนินงาน (semester 1)
% Counts from cyc-monorepo main at commit 5ce62d1 (25 Nov 2026) and tab:progress-status.
\section{การประเมินผลการดำเนินงาน}
\subsection{ความคืบหน้าเทียบกับแผน}
\thaijustify{
  จากงาน 27 รายการใน milestone M0 ถึง M3 (ตารางที่~\ref{tab:progress-status}) งานที่แล้วเสร็จมี 16 รายการ ดำเนินการบางส่วน 10 รายการ และยังไม่เริ่ม 1 รายการ M2 แล้วเสร็จครบทุกรายการ และ M3 ซึ่งเป็นส่วนหลักของภาคการศึกษานี้ใช้งานได้ตั้งแต่การเขียน user story จนถึงการปรับแก้ UI ผ่านแชต งานที่ยังค้างส่วนใหญ่อยู่ใน M1 คือการสมัครสมาชิกด้วยอีเมลและรหัสผ่าน ซึ่งอยู่ระหว่างพัฒนาใน branch แยก และงานที่ต่อเนื่องไปยังการทดสอบ ได้แก่ การมอบหมายงานให้ผู้เข้าร่วมทดสอบและการเริ่มเซสชัน ซึ่งพัฒนาเฉพาะ backend ไว้แล้ว ส่วนงานที่ยังไม่เริ่มคือชุดข้อมูลสำหรับประเมินคุณภาพผลลัพธ์ของ AI
}
\subsection{งานที่ทำได้ล่วงหน้า}
\thaijustify{
  ทีมพัฒนา backend ของ M4 ถึง M6 ไว้ก่อนแผน และออกแบบหน้าจอ high-fidelity ครบทุกกลุ่มผู้ใช้แล้ว รวมถึงหน้าจอของผู้เข้าร่วมทดสอบและผู้มีส่วนได้ส่วนเสียที่จะพัฒนาในภาคการศึกษาที่ 2 (หัวข้อ~\ref{sec:ui-prototype}) จึงเริ่มพัฒนาหน้าจอของภาคการศึกษาที่ 2 ได้ทันทีโดยไม่ต้องรอการออกแบบ
}
\subsection{คุณภาพของซอฟต์แวร์}
\thaijustify{
  ณ commit 5ce62d1 (25 พฤศจิกายน 2569) ใน branch หลัก backend มีไฟล์ unit test 24 ไฟล์ และ end-to-end test 4 ไฟล์ ครอบคลุมการยืนยันตัวตน การควบคุมสิทธิ์ cache และการสร้าง UI การเปลี่ยนแปลงทั้งหมดผ่าน pull request ซึ่ง merge แล้ว 16 ครั้ง และทุก pull request ผ่านการตรวจ lint และ build อัตโนมัติก่อน merge
}
\subsection{คุณภาพผลลัพธ์ของ AI}
\thaijustify{
  ในภาคการศึกษานี้ระบบควบคุมคุณภาพของ UI ที่ AI สร้างด้วยการตรวจอัตโนมัติก่อนแสดงผล ได้แก่ การใช้ component ของ UI kit สีแบรนด์ platform และมาตรฐานการเข้าถึง ร่วมกับการให้โมเดลแก้ไขเมื่อไม่ผ่าน (หัวข้อ~\ref{sec:ai-pipeline}) แต่ยังไม่ได้ประเมินคุณภาพของผลลัพธ์เชิงปริมาณ เช่น สัดส่วนของ variant ที่ผ่านการตรวจในครั้งแรก หรือการให้คะแนนโดยนักออกแบบ ซึ่งจะทำพร้อมกับการจัดทำชุดข้อมูลประเมินในภาคการศึกษาที่ 2
}
